% ============================================================================
% Fichier  : partie4/P04_C19_applications.tex
% Titre    : Sécurité des Applications
% Auteur   : MINKA MI NGUIDJOÏ Thierry Emmanuel
% Version  : 1.0 -- novembre 2026
% ============================================================================
% Base : brouillon « Livre IX » (méthodes de développement, SSDLC,
% vulnérabilités, tests). Sources : NIST SP 800-218 (SSDF, 2022) ; Shostack
% (2014, STRIDE) ; OWASP Top 10 (2021) ; OWASP ASVS (4.0) ; OWASP SAMM.
% Opération MEFOU : portail de paiement fictif de l'ANPP.
% ============================================================================

\chapter{Sécurité des Applications}
\label{chap:applications}
\chaptermark{Sécurité des applications}

\epigraph{%
    « Toute donnée qui vient de l'extérieur est hostile\\
    jusqu'à preuve du contraire. »%
}{Principe du cours}

\niveauChapitre{Intermédiaire}{Programmation (Python ou PHP), bases de
données SQL~; Chapitres~\ref{chap:proprietes} et~\ref{chap:iam}.}

\begin{boxobjectifs}
\begin{itemize}
    \item Intégrer la sécurité à chaque étape du cycle de développement.
    \item Modéliser les menaces d'une application avec STRIDE.
    \item Reconnaître les vulnérabilités de l'OWASP Top~10 et les corriger.
    \item Écrire des requêtes paramétrées et valider les entrées.
    \item Choisir entre revue de code, SAST, DAST et analyse des
          dépendances.
\end{itemize}
\end{boxobjectifs}

\bigskip

% ============================================================================
\section*{Question directrice : comment un numéro dans une adresse a-t-il exposé des milliers de bénéficiaires~?}
% ============================================================================

L'ANPP, agence publique de paiement, ouvre un portail où chaque
bénéficiaire consulte ses paiements à l'adresse
\texttt{/paiements?id=10452}. Un étudiant remarque qu'en changeant le
numéro, il voit les paiements d'un autre. L'application vérifiait que
l'utilisateur était connecté, pas qu'il consultait \emph{ses} données. Ce
chapitre traite de ces erreurs de conception et de code, et de la manière
de les éviter.

\medskip
\textbf{Plan du chapitre.}
\begin{enumerate}
    \item Le cycle de développement sécurisé (section~\ref{sec:c19-ssdlc}).
    \item Modéliser les menaces (section~\ref{sec:c19-stride}).
    \item L'OWASP Top~10 (section~\ref{sec:c19-top10}).
    \item Les injections (section~\ref{sec:c19-injection}).
    \item Le contrôle d'accès applicatif (section~\ref{sec:c19-acces}).
    \item Tester la sécurité d'une application (section~\ref{sec:c19-tests}).
\end{enumerate}

\separateur

% ============================================================================
\section{Le Cycle de Développement Sécurisé}
\label{sec:c19-ssdlc}
% ============================================================================

Le NIST a publié un cadre de développement logiciel sécurisé (SSDF), qui
organise les pratiques en quatre groupes~: préparer l'organisation,
protéger le logiciel, produire un logiciel bien sécurisé, et répondre aux
vulnérabilités~\cite{nist800218}.

\begin{tableaucours}{L{2.8cm}L{9.4cm}}{%
    \tableauHeader{Étape} &
    \tableauHeader{Activité de sécurité}
}
Exigences & Exigences de sécurité (authentification, journalisation,
données sensibles), par exemple tirées de l'ASVS \\
Conception & Modélisation des menaces~; choix d'architecture \\
Développement & Règles de codage~; bibliothèques reconnues~; pas de
secret dans le code~; revue par les pairs \\
Vérification & Analyse statique, tests dynamiques, analyse des
dépendances~; tests d'intrusion pour les applications exposées \\
Mise en production & Configuration durcie~; secrets dans un coffre \\
Maintenance & Correctifs, veille sur les dépendances, traitement des
signalements (Chapitre~\ref{chap:ethique}) \\
\end{tableaucours}
\vspace{-0.8em}
\begin{center}{\small\itshape Tableau~19.1. La sécurité dans le cycle de
développement}\end{center}

% ============================================================================
\section{Modéliser les Menaces}
\label{sec:c19-stride}
% ============================================================================

La \termecle{modélisation des menaces} examine la conception d'une
application pour trouver ce qui peut mal tourner, avant d'écrire le code.
La méthode STRIDE classe les menaces en six catégories, chacune opposée à
une propriété~\cite{shostack2014}.

\begin{tableaucours}{L{3.6cm}L{3.4cm}L{5.2cm}}{%
    \tableauHeader{Menace} &
    \tableauHeader{Propriété violée} &
    \tableauHeader{Exemple sur le portail de l'ANPP}
}
Usurpation (\textit{Spoofing}) & Authenticité & Se connecter avec le
compte d'un autre \\
Altération (\textit{Tampering}) & Intégrité & Modifier le RIB d'un
bénéficiaire \\
Répudiation & Non-répudiation & Nier avoir modifié un RIB, faute de
journal \\
Divulgation (\textit{Information disclosure}) & Confidentialité & Voir les
paiements d'un autre \\
Déni de service & Disponibilité & Saturer la recherche \\
Élévation de privilèges & Autorisation & Accéder aux fonctions
d'administration \\
\end{tableaucours}
\vspace{-0.8em}
\begin{center}{\small\itshape Tableau~19.2. STRIDE appliqué au portail de l'ANPP}\end{center}

% ============================================================================
\section{L'OWASP Top 10}
\label{sec:c19-top10}
% ============================================================================

L'OWASP Top~10 (2021) liste les dix catégories de risques les plus
fréquentes des applications web~\cite{owasp2021top10}. C'est un outil de
sensibilisation~; pour fixer des exigences vérifiables, on utilise
l'ASVS~\cite{owasp2021asvs}.

\begin{tableaucours}{L{1.6cm}L{5.6cm}L{5.0cm}}{%
    \tableauHeader{Rang} &
    \tableauHeader{Catégorie} &
    \tableauHeader{Mesure principale}
}
A01 & Contrôle d'accès défaillant & Vérifier les droits sur chaque objet \\
A02 & Défaillances cryptographiques & TLS, chiffrement, hachage des mots de
passe \\
A03 & Injection & Requêtes paramétrées, encodage des sorties \\
A04 & Conception non sécurisée & Modélisation des menaces \\
A05 & Mauvaise configuration & Durcissement, pas de pages de débogage \\
A06 & Composants vulnérables & Inventaire et mise à jour des dépendances \\
A07 & Identification et authentification & Second facteur, gestion des
sessions \\
A08 & Intégrité des logiciels et des données & Signer les mises à jour,
contrôler la chaîne de production \\
A09 & Journalisation et surveillance & Journaliser les événements de
sécurité \\
A10 & Falsification de requêtes côté serveur & Filtrer les adresses
appelées par le serveur \\
\end{tableaucours}
\vspace{-0.8em}
\begin{center}{\small\itshape Tableau~19.3. L'OWASP Top~10 (2021)}\end{center}

% ============================================================================
\section{Les Injections}
\label{sec:c19-injection}
% ============================================================================

\begin{boxdef}
\termecle{Injection}~: vulnérabilité par laquelle une donnée fournie par
l'utilisateur est interprétée comme une commande (SQL, système, LDAP).
L'\termecle{injection SQL} en est la forme la plus connue.
\end{boxdef}

\begin{boxoutils}{Une requête vulnérable, puis corrigée}
\begin{lstlisting}[language=Python]
import sqlite3
con = sqlite3.connect("anpp.db")

def paiements_vulnerable(nom):
    # la donnee est collee dans la requete : NE PAS FAIRE
    return con.execute(
        "SELECT * FROM paiements WHERE beneficiaire = '" + nom + "'").fetchall()

def paiements_corrige(nom):
    # requete parametree : la donnee reste une donnee
    return con.execute(
        "SELECT * FROM paiements WHERE beneficiaire = ?", (nom,)).fetchall()
\end{lstlisting}
Avec la valeur \texttt{x' OR '1'='1}, la première fonction renvoie tous
les paiements~: la condition devient toujours vraie. La seconde cherche
un bénéficiaire nommé littéralement \texttt{x' OR '1'='1}, et ne trouve
rien.
\end{boxoutils}

\begin{boxalert}
\textbf{Violation typique.} Corriger une injection en supprimant les
apostrophes.

\smallskip
\textbf{Règle~:} une liste de caractères interdits se contourne toujours.
La correction est structurelle~: requête paramétrée, et validation des
entrées selon ce qui est \emph{attendu} (un identifiant est un nombre).
\end{boxalert}

% ============================================================================
\section{Le Contrôle d'Accès Applicatif}
\label{sec:c19-acces}
% ============================================================================

\begin{boxscenario}{Opération MEFOU. Le paramètre id}
\textbf{Constatation.} Le portail de l'ANPP renvoie les paiements de
l'enregistrement dont le numéro figure dans l'adresse, pour tout
utilisateur connecté. Les numéros sont consécutifs.

\medskip
\textbf{Analyse} (confiance élevée). C'est une référence directe à un
objet non sécurisée (catégorie A01)~: l'application vérifie
l'authentification, pas l'autorisation sur l'objet. La correction~:
vérifier côté serveur, à chaque requête, que l'enregistrement appartient
à l'utilisateur connecté, et journaliser les refus. Des numéros
imprévisibles compliquent l'attaque mais ne la corrigent pas.
\end{boxscenario}

% ============================================================================
\section{Tester la Sécurité d'une Application}
\label{sec:c19-tests}
% ============================================================================

\begin{tableaucours}{L{3.0cm}L{4.8cm}L{4.4cm}}{%
    \tableauHeader{Technique} &
    \tableauHeader{Principe} &
    \tableauHeader{Trouve bien}
}
Revue de code & Lecture par un pair, guidée par une liste & Défauts de
logique, contrôle d'accès \\
SAST & Analyse automatique du code source & Injections, secrets dans le
code \\
DAST & Attaque automatique de l'application en fonctionnement &
Configuration, injections exposées \\
Analyse des dépendances & Comparaison des bibliothèques aux CVE &
Composants vulnérables (A06) \\
Test d'intrusion & Attaque manuelle et autorisée & Enchaînements, logique
métier (Chapitre~\ref{chap:evaluer}) \\
\end{tableaucours}
\vspace{-0.8em}
\begin{center}{\small\itshape Tableau~19.4. Techniques de test des applications}\end{center}

\begin{boxmetier}{Développeur}
\textbf{Les réflexes du développeur.} Requêtes paramétrées~; contrôle des
droits sur chaque objet, côté serveur~; aucun secret dans le code ni dans
le dépôt~; dépendances inventoriées et mises à jour~; journalisation des
événements de sécurité~; revue par un pair avant chaque fusion.
\end{boxmetier}

\begin{boxmetier}{Auditeur de sécurité}
\textbf{Auditer un développement.} L'auditeur demande les exigences de
sécurité, la modélisation des menaces, les résultats des analyses
automatiques et leur traitement, et vérifie par sondage que les
vulnérabilités signalées ont été corrigées.
\end{boxmetier}

\separateur

\begin{boxresume}
\textbf{Ce qu'il faut retenir de ce chapitre~:}
\begin{itemize}
    \item \textbf{La sécurité s'intègre au cycle de développement}, des
          exigences à la maintenance (SSDF du NIST).
    \item \textbf{STRIDE} classe les menaces en six catégories, une par
          propriété.
    \item \textbf{L'OWASP Top~10} sensibilise~; l'ASVS fixe des exigences
          vérifiables.
    \item \textbf{L'injection} se corrige par des requêtes paramétrées et
          une validation positive.
    \item \textbf{Le contrôle d'accès} se vérifie sur chaque objet, côté
          serveur.
    \item \textbf{Revue, SAST, DAST, dépendances et intrusion} se
          complètent.
\end{itemize}
\end{boxresume}

\begin{boxinfo}
\textbf{Points de contrôle de l'élève.} Je sais~: placer une activité de
sécurité à chaque étape du cycle~; appliquer STRIDE~; citer les catégories
du Top~10~; corriger une injection SQL~; expliquer une référence directe
non sécurisée~; choisir une technique de test.
\end{boxinfo}

\bigskip

\begin{boxexo}[title={\ding{45}\quad Exercice~19.1. STRIDE \niveaubase}]
Associez une catégorie STRIDE à~: un utilisateur nie avoir passé une
commande~; un attaquant modifie un prix dans le panier~; une page
d'erreur affiche le mot de passe de la base.
\end{boxexo}

\begin{boxexo}[title={\ding{45}\quad Exercice~19.2. Top 10 \niveaubase}]
Classez dans le Top~10~: une bibliothèque de 2017 jamais mise à jour~; des
mots de passe hachés avec MD5~; un mode débogage actif en production.
\end{boxexo}

\begin{boxexo}[title={\ding{45}\quad Exercice~19.3. L'injection \niveaubase}]
Expliquez pourquoi la valeur \texttt{x' OR '1'='1} renvoie tous les
paiements avec la fonction vulnérable.
\end{boxexo}

\begin{boxexo}[title={\ding{45}\quad Exercice~19.4. Corriger \niveaumoyen}]
Réécrivez en requête paramétrée une requête qui recherche les expéditions
d'un client par son numéro et sa ville.
\end{boxexo}

\begin{boxexo}[title={\ding{45}\quad Exercice~19.5. Le paramètre id \niveaumoyen}]
Écrivez, en pseudo-code, la vérification que le portail de l'ANPP aurait
dû faire, et la ligne de journal en cas de refus.
\end{boxexo}

\begin{boxexo}[title={\ding{45}\quad Exercice~19.6. Secrets \niveaumoyen}]
Un développeur a publié sur un dépôt de code le fichier de configuration
contenant le mot de passe de la base. Que faites-vous, dans l'ordre~?
\end{boxexo}

\begin{boxexo}[title={\ding{45}\quad Exercice~19.7. Choisir les tests \niveaumoyen}]
Pour le portail de suivi de LiZbiZ, proposez une combinaison de tests
avant la mise en production, et leur ordre.
\end{boxexo}

\begin{boxexo}[title={\ding{45}\quad Exercice~19.8. Validation positive \niveaumoyen}]
Écrivez les règles de validation des champs~: numéro de suivi (dix
chiffres), adresse électronique, montant en FCFA.
\end{boxexo}

\begin{boxexo}[title={\ding{45}\quad Exercice~19.9. Modéliser \niveauexpert}]
Modélisez avec STRIDE l'application de paiement des fournisseurs de LiZbiZ
(saisie, validation, virement), et proposez une mesure par menace.
\end{boxexo}

\begin{boxexo}[title={\ding{45}\quad Exercice~19.10. Exigences \niveauexpert}]
Rédigez dix exigences de sécurité vérifiables, inspirées de l'ASVS, pour
le cahier des charges de la refonte du portail de l'ANPP.
\end{boxexo}

\bigskip

\begin{boxinfo}
\textbf{Pour aller plus loin.}
\begin{itemize}
    \item NIST SP~800-218 (SSDF)~\cite{nist800218}.
    \item A.~Shostack, \emph{Threat Modeling}~\cite{shostack2014}.
    \item OWASP Top~10 et ASVS~\cite{owasp2021top10,owasp2021asvs}.
\end{itemize}
\end{boxinfo}

\begin{boxinfo}
\textbf{Transition.} De plus en plus d'applications vivent dans le nuage
et sur les téléphones. Le Chapitre~\ref{chap:nuage} en traite les risques
propres.
\end{boxinfo}
